\documentclass[11pt]{article}

\usepackage[margin=1in]{geometry}
\usepackage{amsmath,amssymb,amsfonts}
\usepackage{bm}
\usepackage{graphicx}
\usepackage{booktabs}
\usepackage{enumitem}
\usepackage[numbers,sort&compress]{natbib}
\usepackage[colorlinks=true,linkcolor=blue,citecolor=teal,urlcolor=blue]{hyperref}

\newcommand{\Bvec}{\mathbf{b}}
\newcommand{\nvec}{\mathbf{n}}
\newcommand{\xvec}{\mathbf{x}}
\newcommand{\uvec}{\mathbf{u}}
\newcommand{\tvec}{\mathbf{t}}
\newcommand{\bdry}{\Gamma}
\newcommand{\dom}{\Omega}
\newcommand{\epspl}{\bm{\varepsilon}^{\mathrm{p}}}
\newcommand{\epsstar}{\bm{\varepsilon}^{*}}
\newcommand{\sig}{\bm{\sigma}}
\newcommand{\eps}{\bm{\varepsilon}}
\newcommand{\Cstiff}{\mathbb{C}}
\newcommand{\half}{\tfrac{1}{2}}

\title{\textbf{Modelling Plastic Deformation in VCEM:\\
A Plan for Two Tungsten Fracture Problems}}
\author{N.\ M.\ Ghoniem \\ \small Mechanical \& Aerospace Engineering, UCLA}
\date{\today}

\begin{document}
\maketitle

\begin{abstract}
This note proposes the most efficient routes to introduce plastic deformation
into \texttt{VCEM} (the dislocation-based, boundary-element fracture code in
\texttt{vcem/}; \citealp{Ghoniem2026VCEMa,Ghoniem2026VCEMb}) for two target problems:
(1)~a constrained, thermally shocked disk whose residual plastic strain drives
\emph{brittle} fracture on cool-down, and
(2)~a 2D polycrystal whose grain-boundary (GB) sliding generates internal
stresses, and whose crack tips carry an explicit \emph{local dislocation
distribution} that produces crack-tip shielding --- the mechanism that controls
the ductile-to-brittle transition (DBTT) of tungsten. Bulk crystal plasticity
is deliberately excluded: residual/internal stresses are supplied as a
superimposed external field, and plasticity is resolved only where it matters
mechanically, i.e.\ in the crack-tip process zone. The recommendation exploits
two structural facts of the existing code: cracks are already represented as
\emph{distributions of edge dislocations} solved by a constrained-least-squares
(KKT) coupling to a boundary-only Kelvin BEM, and the crack solver can already
be driven by an arbitrary spatial stress field through
\texttt{AppliedStress.sigma\_func}. A plastic strain is a dislocation density, a
sliding GB is a displacement discontinuity, and a crack-tip plastic zone is a
\emph{continuous distribution of dislocations} in an emergent yielded region ---
all three map onto the same edge-dislocation kernel and the same density solve
VCEM already has, with no discrete-dislocation particle dynamics. The central
technical questions for Problem~2 --- \emph{how to fix the extent of the
crack-tip zone}, and \emph{how to carry a continuous dislocation density on a
slip-plane fan or a local tip grid} --- are addressed in
Sec.~\ref{sec:croacktip}.
\end{abstract}

\tableofcontents

% =====================================================================
\section{What the current code provides}
\label{sec:current}

The relevant facts about \texttt{vcem/}, established by code inspection:

\begin{enumerate}[leftmargin=1.4em]
\item \textbf{Boundary solver.} A 2D direct-collocation elastostatic BEM with
Kelvin kernels $U$ (displacement) and $T$ (traction)
(\texttt{fracture\_utils/Ubem/bem\_solver.py}, \texttt{cpp/src/bem\_solver.cpp}).
It is \emph{boundary-only}: there is currently \emph{no} domain integral and no
body-force / initial-strain term. Interior stresses are recovered by
Somigliana post-processing on a grid (\texttt{stress\_on\_grid}).

\item \textbf{Crack solver.} Cracks are \emph{distributions of edge
dislocations} (a distributed-dislocation / DDM model,
\citealp{hills1996ddt,bilby1968}). The unknowns are the differential
Burgers-vector densities $d\Bvec(s)$ along each crack; the dislocation stress
and displacement kernels live in \texttt{cpp/src/edge\_dislocation.cpp}. A
constrained least-squares (``KKT'') solve
(\texttt{cpp/src/kkt.cpp}, \texttt{crack\_assemblers.cpp}) couples the crack
densities to the boundary BEM unknowns and enforces crack-opening-displacement
(COD) continuity at junctions.

\item \textbf{Driving field.} \texttt{AppliedStress} already accepts a spatial
field \texttt{sigma\_func}$(\xvec)\!\to\!\sig(\xvec)$ and feeds it to the crack
solve. The docstring states this ``enables driving an infinite-medium crack
solve with a BEM-computed stress field without coupling.''

\item \textbf{SIF and propagation.} $(K_I,K_{II})$ are extracted from a
Peach--Koehler window force (\texttt{Uprocessor/SIF\_pk.py}); the kink direction
uses the maximum-hoop-stress / Erdogan--Sih criterion
\citep{erdogan_sih1963}; growth is quasi-static with a toughness interface
(\texttt{Upropagation/toughness.py}) that already supports a \emph{spatially
varying} $K_c(\xvec)$ via \texttt{CallableToughness}.

\item \textbf{Material model.} \texttt{Material(E, nu, plane\_stress)} is linear
isotropic elasticity only. There is no plasticity, no eigenstrain, no thermal
strain, and no notion of grains anywhere in the tree.
\end{enumerate}

\noindent
The two structural assets in items 2--3 are decisive: a plastic strain enters
linear elasticity exactly as an \emph{eigenstrain} \citep{mura1987,eshelby1957},
and an eigenstrain field is kinematically equivalent to a dislocation density.
Because VCEM is already a superposition engine over dislocation sources plus a
boundary-correction solve, we can add plasticity \emph{without} building a full
domain-integral BEM in most cases.

% =====================================================================
\section{The unifying principle: plasticity as eigenstrain}
\label{sec:eigenstrain}

Decompose the total strain into elastic and inelastic parts,
$\eps = \eps^{\mathrm e} + \epsstar$, where the eigenstrain $\epsstar$ collects
all stress-free inelastic contributions:
\begin{equation}
\epsstar \;=\; \epspl \;+\; \eps^{\mathrm{th}},
\qquad
\eps^{\mathrm{th}} = \alpha\,\Delta T(\xvec)\,\mathbf{I},
\label{eq:eigen-decomp}
\end{equation}
with $\epspl$ the plastic strain and $\eps^{\mathrm{th}}$ the thermal strain.
The constitutive law is
\begin{equation}
\sig = \Cstiff : (\eps - \epsstar) = \Cstiff : \eps^{\mathrm e}.
\label{eq:hooke-eigen}
\end{equation}
Mura's equivalence \citep{mura1987} states that the stress field of an
eigenstrain $\epsstar$ in a body equals that of an equivalent body-force
distribution $\mathbf f = -\nabla\!\cdot(\Cstiff:\epsstar)$, or equivalently a
dislocation density
\begin{equation}
\bm{\alpha}_{\mathrm{disl}} = -\,\nabla\times\bm{\beta}^{\mathrm p},
\qquad \bm{\beta}^{\mathrm p}\ \text{the plastic distortion},\ \
\epspl=\half(\bm{\beta}^{\mathrm p}+\bm{\beta}^{\mathrm p\,\mathsf T}).
\label{eq:disl-density}
\end{equation}

\paragraph{Finite-body residual stress by superposition (Bueckner/Erdogan).}
For a finite domain $\dom$ with boundary $\bdry$, the residual stress from
$\epsstar$ is obtained in two steps that reuse VCEM as-is:
\begin{equation}
\sig(\xvec) \;=\; \underbrace{\sig^{\infty}_{\epsstar}(\xvec)}_{\text{infinite-medium eigenstrain field}}
\;+\;
\underbrace{\sig^{\mathrm{corr}}(\xvec)}_{\text{BEM correction}} ,
\label{eq:superpose}
\end{equation}
where $\sig^{\mathrm{corr}}$ is the elastic field that cancels the spurious
tractions $\tvec^{\infty}_{\epsstar}=\sig^{\infty}_{\epsstar}\nvec$ left on the
true boundary $\bdry$, i.e.\ the BEM is solved with boundary condition
$\tvec^{\mathrm{corr}} = -\,\tvec^{\infty}_{\epsstar}$ on traction segments (and
the matching displacement correction on displacement segments). This is exactly
the superposition VCEM already uses to seat dislocation-based cracks inside a
finite boundary; the eigenstrain simply contributes another set of source
tractions. \emph{No new domain integral on the boundary solver is needed} as
long as we can evaluate $\sig^{\infty}_{\epsstar}$ at boundary collocation and
field points.

\paragraph{When a domain integral is unavoidable.} If $\epsstar$ is genuinely
volumetric and spread over a large fraction of $\dom$ (full-field bulk CP), the
classical initial-strain BEM \citep{telles1983,brebbia1984,aliabadi2002} adds a
domain term to the boundary integral equation,
\begin{equation}
c_{ij}u_j(\mathbf p)+\!\int_{\bdry}\!T_{ij}u_j\,d\bdry
=\!\int_{\bdry}\!U_{ij}t_j\,d\bdry
+\int_{\dom} \! E_{ijk}\,\epsstar_{jk}\,d\dom,
\label{eq:bie-domain}
\end{equation}
requiring volume cells and quadrature. This is the expensive route and we use it
only where Sec.~\ref{sec:problem2} truly needs full-field plasticity.

% =====================================================================
\section{Problem 1 --- Thermal-shock residual strain, brittle fracture}
\label{sec:problem1}

\paragraph{Physics.} A constrained disk is heated (thermal shock), yields in
compression where it is hottest/most constrained, then on cooling the locked-in
plastic strain leaves a self-equilibrated \emph{residual} stress field
$\sig^{\mathrm{res}}$. Tensile regions of $\sig^{\mathrm{res}}$ drive
\emph{brittle} cracking. Because the material is brittle, crack growth does
\emph{not} feed back into the plastic flow --- the coupling is strictly one-way,
plasticity $\rightarrow$ stress $\rightarrow$ fracture.

\paragraph{Recommendation: one-way eigenstrain superposition (Option 1, done
correctly).} For a one-way brittle problem this is not a shortcut --- it is
exact. The plan:
\begin{enumerate}[leftmargin=1.4em]
\item Obtain the residual eigenstrain field $\epsstar(\xvec)=\epspl+\alpha\Delta T$.
  Two interchangeable sources:
  \begin{itemize}[leftmargin=1.2em]
  \item \textbf{(1a) Import from COMSOL.} Run the thermo-elasto-plastic
    transient of the constrained disk in COMSOL, export $\sig^{\mathrm{res}}(\xvec)$
    (or $\epspl(\xvec)$) on a mesh, and wrap a SciPy interpolant as
    \texttt{sigma\_func}. Lowest effort; leverages a validated FE plasticity
    solver. The exported field already contains the finite-body boundary
    reactions, so Eq.~\eqref{eq:superpose} reduces to a direct hand-off.
  \item \textbf{(1b) In-house thermal-shock model.} A self-contained transient
    thermo-elasto-plastic solver built on VCEM's own dislocation/initial-strain
    machinery, removing the COMSOL dependency. Developed in full in
    Sec.~\ref{sec:inhouse}.
  \end{itemize}
\item Form the driving field by Eq.~\eqref{eq:superpose}: if using (1b) on an
  infinite-medium basis, let VCEM's BEM compute $\sig^{\mathrm{corr}}$ to satisfy
  the disk's true (traction-free or constrained) boundary; if using (1a) on the
  actual disk geometry, pass $\sig^{\mathrm{res}}$ straight through.
\item Run the existing brittle propagation: feed $\sig^{\mathrm{res}}$ via
  \texttt{AppliedStress.sigma\_func}; nucleate/extend cracks with the existing
  $(K_I,K_{II})$ + Erdogan--Sih + $K_c$ machinery. Tungsten cleavage anisotropy,
  if desired, enters through a spatially/orientationally varying $K_c(\xvec)$ in
  \texttt{CallableToughness}.
\end{enumerate}

\paragraph{Effort and risk.} Low. Path (1a) is essentially a config/IO task on
top of capabilities the code already exposes; path (1b) adds one self-contained
$\sim$200--400 line plasticity module with no change to the BEM/KKT core. The
only correctness subtlety is ensuring the residual field is
\emph{self-equilibrated} on the true boundary --- handled automatically by
either COMSOL's actual-geometry solve or by the
Eq.~\eqref{eq:superpose} BEM correction. This is the recommended starting point
and also a clean validation vehicle (it quantifies how large the
plasticity-induced correction to the SIFs actually is) before investing in
Problem~2.

\subsection{In-house thermal-shock model (self-contained alternative)}
\label{sec:inhouse}

When a COMSOL hand-off is undesirable, the residual field can be produced
entirely inside VCEM with three coupled ingredients. The model is transient and
one-way to fracture, so it is solved \emph{once} per loading history and its
output is the residual stress field of Sec.~\ref{sec:problem1}.

\paragraph{(i) Transient heat conduction.} Solve the heat equation for the
thermal-shock pulse (e.g.\ a surface heat flux $q''(t)$ representing a plasma
transient on a tungsten face) \citep{carslaw_jaeger1959}:
\begin{equation}
\rho\,c_p(T)\,\frac{\partial T}{\partial t}
= \nabla\!\cdot\!\big(k(T)\,\nabla T\big),
\qquad
-k\,\frac{\partial T}{\partial n}\Big|_{\bdry_q}=q''(t),
\label{eq:heat}
\end{equation}
on a background mesh covering the disk. For axisymmetric or mildly 2D shocks a
short finite-difference/finite-element thermal solve is adequate; the output is
$T(\xvec,t)$ and hence the thermal eigenstrain
$\eps^{\mathrm{th}}=\alpha(T)\,\Delta T(\xvec,t)\,\mathbf I$.

\paragraph{(ii) Incremental thermo-elasto-plastic solve via initial-strain BEM.}
This is the one place a domain term is genuinely warranted, because
thermal-shock plasticity is volumetric. Tessellate the disk interior into
\emph{internal cells} and use the initial-strain boundary integral equation
\citep{telles1983,brebbia1984,aliabadi2002}, with both the thermal and plastic
strains carried as the eigenstrain in Eq.~\eqref{eq:bie-domain}. At each time
increment, the stress at cell integration points drives a von Mises
($J_2$) return mapping with temperature-dependent yield
\citep{simo_taylor1985}:
\begin{equation}
f=\sqrt{\tfrac{3}{2}\,\mathbf s\!:\!\mathbf s}-\sigma_Y(T,\bar\varepsilon^{\mathrm p})\le 0,
\qquad
\Delta\epspl=\Delta\lambda\,\frac{\partial f}{\partial \sig}
=\Delta\lambda\,\sqrt{\tfrac{3}{2}}\,\frac{\mathbf s}{\|\mathbf s\|},
\label{eq:j2-inhouse}
\end{equation}
iterated to satisfy consistency ($f\le0$, $\Delta\lambda\,f=0$) over the cells
each step. The same edge-dislocation kernels already in VCEM supply the cell
stress influence coefficients --- the initial-strain BEM \emph{is} a dislocation
(eigenstrain) superposition, so the C++ kernel layer is reused; only 2D cell
quadrature and the return-mapping loop are new.

\paragraph{(iii) Residual field and hand-off.} After the cool-down increment the
locked-in $\epspl(\xvec)$ defines the residual eigenstrain; its finite-body
stress (boundary tractions cancelled by the BEM correction,
Eq.~\eqref{eq:superpose}) is the self-equilibrated $\sig^{\mathrm{res}}(\xvec)$
passed to the brittle solver exactly as in path~(1a). Validation target: a
constrained-disk thermal shock should reproduce the known sign pattern
(compressive yield at the hot face, tensile residual zone on cool-down) and
crack onset consistent with W heat-load experiments \citep{araksheev_w,
pintsuk2011}.

\paragraph{Effort.} Medium. It reuses the dislocation kernels but adds (a) a
transient thermal solver, (b) internal-cell domain quadrature, and (c) the
return-mapping loop. Note the internal-cell infrastructure built here is
\emph{exactly} what the optional 2D area-distribution route in
Sec.~\ref{sec:croacktip} would need, so the investment is shared with
Problem~2.

% =====================================================================
\section{Problem 2 --- Polycrystal, GB sliding, and crack-tip plasticity (DBTT)}
\label{sec:problem2}

\paragraph{Physics and scope.} A 2D aggregate of grains is loaded externally.
GB sliding accommodates intergranular shear and builds stress concentrations at
triple junctions and GB ledges; any volumetric residual/internal stress is
supplied as a \emph{superimposed external field} (Sec.~\ref{sec:problem1}), so
\textbf{bulk crystal plasticity is excluded by design}. Plasticity is resolved
\emph{only} where it controls the fracture outcome --- in the process zone of
each active crack tip, as an explicit cloud of dislocations on slip planes. The
DBTT then emerges from the competition between cleavage (lattice or GB) and the
crack-tip shielding that this dislocation cloud provides, with temperature
entering through dislocation \emph{mobility} and \emph{nucleation}
\citep{rice_thomson1974,rice1992,hirsch_roberts1991,gumbsch1998,
giannattasio_roberts2007,tarleton_roberts2009,hartmaier_gumbsch2005}.

\subsection{Microstructure and grain boundaries}
Tessellate the disk into grains by a Voronoi/Laguerre construction
\citep{voronoi_okabe2000}, producing a planar graph of GB segments and triple
junctions --- a structure the \texttt{Ugenerator} topology layer is well suited
to host. Each grain carries an orientation $\theta_g$ fixing its in-plane slip
geometry (BCC tungsten $\{110\}/\{112\}\langle111\rangle$ projected to 2D as
slip directions $\mathbf s^{\alpha}$ and normals $\mathbf m^{\alpha}$). A GB
segment is a displacement-discontinuity interface --- \emph{the same object
VCEM already discretizes for cracks} --- endowed with a cohesive/friction
traction--separation law \citep{xu_needleman1994,raj_ashby1971}:
\begin{equation}
t_s = \eta_{\mathrm{gb}}(T)\,\dot\delta_s \ \ (\text{sliding}),
\qquad
\text{decohere when } t_n \ge \sigma_{\mathrm{gb}}^{\max}(T),
\label{eq:gb}
\end{equation}
so the brittle intergranular branch is captured purely with discontinuity
elements already native to the solver.

% --------------------------------------------------------------
\subsection{Crack-tip plasticity as a continuous dislocation distribution}
\label{sec:croacktip}

Discrete dislocation dynamics is the wrong tool to graft onto VCEM, for two
reasons: (i)~the lattice Burgers vector is so small that a realistic tip zone
would contain an enormous number of individual dislocations, and (ii)~tracking
their glide is an ODE time-integration framework alien to VCEM's
density/integral-equation solve. We therefore represent the tip plastic zone as
a \emph{continuous distribution of dislocations} --- the \emph{same kind of
object already used for crack faces}, where a smooth density carries the Burgers
content and $b$ never appears as a counted quantum
\citep{bilby1968,hills1996ddt,bcs1963,weertman1996,lardner1974}. The zone is an
\emph{emergent yielded region} whose density is solved, per load increment, from
a (singular) integral equation under a flow condition --- a sequence of linear
solves in exactly VCEM's framework, not a particle dynamics. The two questions
this raises are \emph{(Q1)} how to fix the extent and \emph{(Q2)} which
geometric carrier to use.

\subsubsection{Q2 --- carrier: slip-plane fan or local tip grid}
Both options keep the solve in VCEM's wheelhouse.

\paragraph{(V1) Continuous density on a slip-plane fan (recommended for DBTT).}
Emanate a small fan of slip planes (rays) from the tip along the grain's slip
directions $\mathbf s^{\alpha}$. On plane $\alpha$ carry a continuous glide
density $\rho_{\alpha}(r)=-\,d\delta_{\alpha}/dr$, with $\delta_{\alpha}(r)$ the
plastic slip and $r$ the distance from the tip --- identical in kind to the
crack-opening density VCEM already solves on a polyline. The governing condition
replaces the crack's traction-free face by a \emph{flow (yield) condition} on the
yielded part $0\le r\le\omega_{\alpha}$:
\begin{equation}
\tau^{\mathrm{ext}}_{\alpha}(r)
+\sum_{\alpha'}\int_{0}^{\omega_{\alpha'}}\!\!
K_{\alpha\alpha'}(r,r')\,\rho_{\alpha'}(r')\,dr'
=\tau_f(T,\dot\varepsilon)\,\operatorname{sgn}(\rho_\alpha),
\qquad |\tau|<\tau_f \ \text{elsewhere},
\label{eq:cdd}
\end{equation}
where $\tau^{\mathrm{ext}}_{\alpha}$ is the resolved shear from the crack,
applied/residual field, GBs and the BEM boundary correction, $K_{\alpha\alpha'}$
is the (self and cross-plane) shear-influence kernel built from the
\emph{existing} edge-dislocation kernel (\texttt{cpp/src/edge\_dislocation.cpp}),
and $\tau_f$ is the lattice friction/flow stress. This is the planar
continuously-distributed-dislocation generalization of the
Bilby--Cottrell--Swinden strip-yield model \citep{bcs1963,chang_ohr1981}. Because
it carries the crystallographic slip direction explicitly, V1 is the natural
choice for the orientation-dependent polycrystal DBTT, and its assembler is the
crack-face solve with a different right-hand side.

\paragraph{(V2) Continuous density on a local tip grid.} If a smooth,
orientation-averaged zone is preferred, attach a small 2D grid of cells to the
tip and carry a continuous plastic strain $\epspl(\xvec)$ (eigenstrain) on it.
The stress is the initial-strain BEM field of Eq.~\eqref{eq:bie-domain} ---
\emph{the same machinery as the in-house thermal-shock solve}
(Sec.~\ref{sec:inhouse}) --- and the density follows from a $J_2$ flow condition
(Eq.~\eqref{eq:j2-inhouse}) with $\sigma_Y(T)$. Heavier than V1 (genuine 2D cell
quadrature, near-singular self-cell treatment --- cf.\ the interior-eval
near-singularity already documented in the BEM) but constitutively simpler and a
direct reuse of Sec.~\ref{sec:inhouse}. Use V2 when the slip-system geometry is
not the point; use V1 when it is.

\subsubsection{Fan orientation and count (the key V1 modelling choice)}
\label{sec:fan}
Two ways to orient the slip planes, with materially different anisotropy:

\paragraph{(a) Crystallographic ($\theta_g$-locked) --- recommended.} Place the
planes along the in-plane traces of the BCC tungsten slip systems
$\{110\}/\{112\}\langle111\rangle$, rotated by the grain orientation $\theta_g$.
A plane carries density only when its Schmid-resolved shear reaches flow,
\begin{equation}
\tau^{\alpha}=\sig:\big(\mathbf s^{\alpha}\!\otimes\!\mathbf m^{\alpha}\big),
\qquad
\text{active if } |\tau^{\alpha}|\ge \tau_f(T,\dot\varepsilon),
\label{eq:schmid}
\end{equation}
so the \emph{active subset} of the fan is itself emergent (an active set over
planes, on top of the active set over yielded length). This is the only choice
that reproduces grain-to-grain variability --- why one grain shields and blunts
while its misoriented neighbour cleaves --- which is precisely the polycrystal
DBTT physics. It ties $K_{\mathrm{shield}}$ (Eq.~\eqref{eq:shield}) to the
Schmid factors and makes the predicted transition orientation-dependent.

\paragraph{(b) Maximum-shear (continuum) --- baseline only.} Use one or two
planes along the local maximum-shear directions ($\approx\pm45^\circ$ to the tip
principal stress). Simple, robust, recovers a BCS/Dugdale-like isotropic zone and
is a good check against V2 --- but it discards $\theta_g$ and therefore
\emph{cannot} produce orientation-dependent DBTT or grain variability. Its plane
set also rotates with the evolving mixed-mode field and must be re-aimed each
increment. Recommended only as a validation baseline.

\paragraph{Count and a 2D caveat.} A practical fan is the set of distinct
in-plane $\langle111\rangle$ traces for the chosen 2D section (typically
$3$--$6$ independent directions); more planes add density-unknown blocks but
sharpen the anisotropy. Allow each plane to yield independently --- mixed-mode
tip fields from GB sliding break fan symmetry. \textbf{Caveat:} VCEM's 2D
plane-strain edge-dislocation kernel carries only the \emph{in-plane (edge)}
Burgers content; BCC $\langle111\rangle$ slip has strong screw character whose
out-of-plane (mode-III) shielding is \emph{not} represented. The model therefore
captures edge-component shielding only --- adequate for a plane-strain DBTT study
but a known limitation to state explicitly. A \emph{recommended middle ground}:
crystallographic orientations (a) with the Schmid active-set
Eq.~\eqref{eq:schmid}, retaining only the in-plane edge projection.

\subsubsection{Q1 --- the emergent extent}
The zone size is \emph{not} prescribed. The small-scale-yielding estimate
\begin{equation}
\omega \sim \frac{1}{\beta\pi}\left(\frac{K_I}{\tau_f(T)}\right)^{2}
\ \ (\text{or }\sigma_Y),\qquad \beta=2\,(\sigma)\ \text{or}\ 6\,(\varepsilon),
\label{eq:rp}
\end{equation}
\citep{irwin1957} serves only as the \emph{initial guess} for the active-set
iteration. The actual extent comes from a closure condition:
\begin{itemize}[leftmargin=1.4em]
\item \textbf{V1:} the BCS/Dugdale \emph{bounded-stress closure}
\citep{bcs1963,dugdale1960} --- the yielded length $\omega_{\alpha}$ on each plane
is the value for which the density vanishes smoothly at the elastic--plastic
boundary, i.e.\ the $1/\sqrt{r}$ stress singularity at $r=\omega_{\alpha}$ is
removed. Each increment is a coupled singular-integral solve for $\{\rho_\alpha\}$
with an outer iteration on $\{\omega_\alpha\}$ enforcing closure.
\item \textbf{V2:} the \emph{von Mises yield locus}
$\sigma_{\mathrm{eq}}=\sigma_Y(T)$ --- an active-set sweep grows the plastic cells
until the elastic stress just outside is sub-yield and consistency
($f\le0,\ \Delta\lambda\,f=0$) holds inside.
\end{itemize}
In both cases the temperature-dependent flow stress makes the emergent zone
shrink as $T$ falls --- the seed of the DBTT.

\subsection{Ductile-to-brittle transition criterion}
The continuous density shields the tip; the shielding is an \emph{integral} over
the density (no sum over discrete dislocations):
\begin{equation}
K_{\mathrm{tip}}=K_{\mathrm{applied}}-K_{\mathrm{shield}},
\qquad
K_{\mathrm{shield}}=\sum_{\alpha}\int_{0}^{\omega_{\alpha}}
\frac{\mu}{(\kappa+1)\sqrt{2\pi r}}\;\Phi(\varphi_\alpha)\,\rho_{\alpha}(r)\,dr ,
\label{eq:shield}
\end{equation}
with $\Phi(\varphi_\alpha)$ the angular weight of a glide dislocation on a plane
at angle $\varphi_\alpha$ to the crack \citep{rice_thomson1974,weertman1996}. The
flow stress is thermally activated (BCC kink-pair), giving both temperature and
rate dependence:
\begin{equation}
\tau_f(T,\dot\varepsilon)=\tau_0+\Delta\tau\!\left[1-\!\left(\frac{k_BT}{\Delta H_0}
\ln\frac{\dot\varepsilon_0}{\dot\varepsilon}\right)^{1/q}\right]^{1/p}.
\label{eq:tauf}
\end{equation}
Fracture is brittle when the \emph{local} tip intensity reaches cleavage before
the zone can shield it:
\begin{equation}
\text{cleavage if } K_{\mathrm{tip}}\ge K_{Ic}^{\mathrm{cleave}}
\ \ (\text{lattice or } K_c^{\mathrm{gb}}),
\qquad
\text{else blunting/ductile.}
\label{eq:dbtt}
\end{equation}
At low $T$, $\tau_f$ is large $\Rightarrow$ $\omega_{\alpha}\!\to\!0$,
$K_{\mathrm{shield}}\!\to\!0$, and cleavage/GB fracture wins; at high $T$ the
yielded zone extends, $K_{\mathrm{shield}}$ grows, and the tip blunts. The
crossover $T$ depends on the loading rate $\dot K$ through
Eq.~\eqref{eq:tauf} --- reproducing the observed rate-dependent W DBTT
\citep{giannattasio_roberts2007,gumbsch1998,hirsch_roberts1991,
hartmaier_gumbsch2005}. This continuum picture also connects to the group's
distributed-dislocation methodology \citep{ghoniem2000parametric,ghoniem_affine}.
The existing spatially varying $K_c(\xvec)$ (\texttt{CallableToughness}) supplies
the weak-GB cleavage field for the intergranular branch.

\subsection{Coupling architecture and phasing}
At each load/temperature increment, run a staggered active-set iteration on the
existing KKT core:
\begin{enumerate}[leftmargin=1.4em,itemsep=2pt]
\item \textbf{Elastic solve.} Assemble the superposed stress
(Eq.~\eqref{eq:superpose}) from crack-face densities, GB discontinuities, the
current plastic-zone density ($\rho_{\alpha}$ in V1 / $\epspl$ in V2) and the
superimposed external/residual field --- BEM enforcing the true outer boundary.
This is the existing solve with extra source rows.
\item \textbf{Plastic solve.} Update the zone density and length from the flow
condition (Eq.~\eqref{eq:cdd} + closure for V1; return mapping + active set for
V2). \emph{Linear density solves, no time integration.}
\item \textbf{GB update.} Update sliding/decohesion from Eq.~\eqref{eq:gb}.
\item \textbf{Fracture check.} Cleave (crack/GB advance) when
Eq.~\eqref{eq:dbtt} holds; otherwise continue blunting.
\item Iterate 1--4 to convergence, then advance the increment.
\end{enumerate}

Because crack faces, GBs, and the tip plastic density are all
continuous-dislocation objects on the \emph{same} kernel, every step shares one
assembler --- far cheaper than a generic CP--FEM/CP--BEM coupling. Phasing:
\textbf{(2A)} elastic grains + GB sliding + brittle decohesion (no tip
plasticity) to validate the discontinuity machinery; \textbf{(2B)} add the V1
continuous tip density and the shielding criterion Eq.~\eqref{eq:shield};
\textbf{(2C)} sweep $T$ and $\dot K$ through Eqs.~\eqref{eq:tauf} and
\eqref{eq:gb} to map the DBTT.

% =====================================================================
\section{Summary of recommendations}
\label{sec:summary}

\begin{table}[h]
\centering
\small
\begin{tabular}{@{}p{2.3cm}p{4.7cm}p{4.7cm}p{2.5cm}@{}}
\toprule
 & \textbf{Problem 1 (thermal shock)} & \textbf{Problem 2 (polycrystal DBTT)} & \textbf{Effort} \\
\midrule
\textbf{Recommended} &
One-way eigenstrain superposition: residual $\sig^{\mathrm{res}}$ from COMSOL
(1a) or the in-house thermal-shock model (1b, Sec.~\ref{sec:inhouse})
$\to$ \texttt{sigma\_func} $\to$ brittle solve &
GBs as discontinuities + crack-tip plasticity as a \emph{continuous dislocation
density} on a slip-plane fan (V1), zone extent \emph{emergent} from BCS closure;
staggered active-set update on the KKT core & 1a: low \newline 1b: medium
\newline 2: high \\[2pt]
\textbf{Why} &
Exact for one-way brittle coupling; hook already exists &
Crack faces, GBs and the tip density are all continuous-dislocation objects $\to$
one assembler, one density solve; DBTT enters through $\tau_f(T,\dot\varepsilon)$
at the tip & --- \\[2pt]
\textbf{Avoid} &
Building a domain-integral BEM (use it only inside 1b) &
Discrete dislocation dynamics (enormous $N$, alien framework); bulk crystal
plasticity & --- \\
\bottomrule
\end{tabular}
\caption{Recommended modelling routes. Start with Problem~1 (1a) as a
fast-validating capability; build Problem~2 in the phased order
2A$\to$2B$\to$2C with the V1 continuous tip density as the workhorse.}
\end{table}

\paragraph{Single most important design decision.} Keep plasticity as a
\emph{continuous dislocation density} solved on the same kernel as crack faces.
Cracks, grain boundaries, and the crack-tip plastic zone are all
continuous-dislocation-content objects expressible through one set of kernels and
one boundary-correction solve, so build crack-tip plasticity as a continuous
density on a slip-plane fan (V1) and let the zone extent emerge from the
BCS bounded-stress closure, with $\tau_f(T,\dot\varepsilon)$ carrying the DBTT.
Reserve the domain integral Eq.~\eqref{eq:bie-domain} for the cases that truly
need a volumetric field --- the in-house thermal-shock solve
(Sec.~\ref{sec:inhouse}) and the optional V2 tip grid. Do \emph{not} adopt
discrete dislocation dynamics (the lattice $b$ forces an intractable count and
its particle-ODE solve is alien to VCEM), and do not introduce bulk crystal
plasticity: the DBTT physics lives at the tip.

\bibliographystyle{unsrtnat}
\bibliography{plasticity_plan}

\end{document}
